\section{Curvature induced by a spherical bulge}
\label{app:curvature-bulge}

To make the connection between extrinsic bending and intrinsic scalar
curvature more explicit, consider a localized spherically symmetric
deformation of a planar brane. Let the undeformed brane lie in the
$(x,y,z)$-hyperplane of $\mathbb{R}^4$, and suppose a Gaussian bump
is added in the fourth direction,
\begin{equation}
X^4(x,y,z) \;=\; A\,\exp\!\Bigl(-\frac{x^2+y^2+z^2}{2\sigma^2}\Bigr),
\label{eq:gaussian-bulge}
\end{equation}
with amplitude $A$ and width $\sigma$. For small amplitude
$A\ll\sigma$, the induced metric on the brane is approximately
\begin{equation}
g_{ij} \;=\; \delta_{ij} \;+\; \partial_i X^4\,\partial_j X^4
\;\approx\; \delta_{ij} \;+\; \mathcal{O}\bigl((A/\sigma)^2\bigr).
\end{equation}
The second fundamental form (extrinsic curvature) is, to leading
order in $A/\sigma$,
\begin{equation}
b_{ij} \;=\; -\,\partial_i\partial_j X^4
\;=\; \frac{A}{\sigma^2}\,\exp\!\Bigl(-\frac{r^2}{2\sigma^2}\Bigr)
\Bigl(\delta_{ij} \;-\; \frac{x_i x_j}{\sigma^2}\Bigr),
\end{equation}
where $r^2=x^2+y^2+z^2$. The mean curvature (trace) is
\begin{equation}
K \;=\; \tr(b)
\;=\; \frac{A}{\sigma^2}\,\exp\!\Bigl(-\frac{r^2}{2\sigma^2}\Bigr)
\Bigl(3 - \frac{r^2}{\sigma^2}\Bigr),
\end{equation}
and the squared norm is
\begin{equation}
|b|^2 \;=\; \frac{A^2}{\sigma^4}\,\exp\!\Bigl(-\frac{r^2}{\sigma^2}\Bigr)
\Bigl(3 + \frac{r^4}{\sigma^4} - 2\,\frac{r^2}{\sigma^2}\Bigr)
\;+\; \mathcal{O}\bigl((A/\sigma)^3\bigr).
\end{equation}
By the Gauss--Codazzi relation~\eqref{eq:gauss-codazzi-3d}, the
three-dimensional scalar curvature is
\begin{equation}
R_3 \;=\; K^2 \;-\; |b|^2
\;=\; \frac{A^2}{\sigma^4}\,\exp\!\Bigl(-\frac{r^2}{\sigma^2}\Bigr)
\Bigl[
\Bigl(3 - \frac{r^2}{\sigma^2}\Bigr)^2
\;-\; \Bigl(3 + \frac{r^4}{\sigma^4} - 2\,\frac{r^2}{\sigma^2}\Bigr)
\Bigr]
\;+\; \mathcal{O}\bigl((A/\sigma)^3\bigr).
\end{equation}
At the center ($r=0$), one obtains
\begin{equation}
R_3(0) \;=\; \frac{6\,A^2}{\sigma^4},
\end{equation}
confirming that a localized out-of-brane bulge of characteristic size
$\sigma$ and amplitude $A$ induces intrinsic curvature of order
$A^2/\sigma^4$ in the brane. This quadratic scaling in amplitude
reflects the geometric origin of $R_3$ as the difference between the
square of the mean curvature and the squared Frobenius norm of the
second fundamental form.